\chapter{Spectral-sequence machinery}
\label{chap:spectral-sequence-machinery}

This chapter packages spectral objects, filtered-complex realizations, page presentations, reindexing, convergence, and comparison interfaces.  It replaces the former catch-all Background chapter.

\section{Mathlib spectral-sequence and spectral-object interfaces}

\begin{definition}[Spectral sequence in an abelian category]
  \label{def:mathlib-spectral-sequence}
  \lean{CategoryTheory.SpectralSequence}
  \mathlibok
  \uses{def:mathlib-abelian,def:mathlib-homological-complex}
  Let $\mathcal C$ be an abelian category, let $\kappa$ be an index type, let
  $c : \mathbb Z \to \operatorname{ComplexShape}(\kappa)$ specify the shape
  of the differential on each page, and fix a starting page $r_0\in\mathbb Z$.
  A spectral sequence consists of a homological complex $E_r$ of shape $c(r)$
  for every $r\ge r_0$, together with, for every $p\in\kappa$, an isomorphism
  \[
    H(E_r)_p \cong (E_{r+1})_p.
  \]
  This is Mathlib's definition \texttt{CategoryTheory.SpectralSequence} in
  the pinned revision.  All generic KIP126 statements use this definition
  directly rather than a project-local synonym.
\end{definition}

% SOURCE: Mathlib/Algebra/Homology/SpectralSequence/Basic.lean, declaration CategoryTheory.SpectralSequence.Hom.
\begin{definition}[Mathlib spectral-sequence morphism]
  \label{def:mathlib-spectral-sequence-hom}
  \lean{CategoryTheory.SpectralSequence.Hom}
  \mathlibok
  A Mathlib morphism between spectral sequences has pagewise chain maps
  compatible with passage to homology, but both sequences must have the same
  index type, complex shape, and starting page.  It therefore cannot express
  the classical--synthetic reindexing required later.
\end{definition}

% SOURCE: Mathlib/CategoryTheory/Triangulated/SpectralObject.lean:49, declaration CategoryTheory.Triangulated.SpectralObject.
\begin{definition}[Mathlib triangulated spectral object]
  \label{def:mathlib-triangulated-spectral-object}
  \lean{CategoryTheory.Triangulated.SpectralObject}
  \mathlibok
  \uses{def:mathlib-shift,def:mathlib-triangle,def:mathlib-pretriangulated}
  For a pretriangulated category $\mathcal C$ and an index category $I$, a
  triangulated spectral object consists of a functor from length-one composable
  arrows in $I$ to $\mathcal C$ and a functorial distinguished triangle for
  every composable pair $i\to j\to k$.  It is an object of
  \texttt{CategoryTheory.Triangulated.SpectralObject}; it is not the abelian
  spectral object in the next node.
\end{definition}

% SOURCE: Mathlib/CategoryTheory/Triangulated/SpectralObject.lean:103, declaration CategoryTheory.Triangulated.SpectralObject.precomp.
\begin{definition}[Precomposition of a triangulated spectral object]
  \label{def:mathlib-triangulated-spectral-object-precomp}
  \lean{CategoryTheory.Triangulated.SpectralObject.precomp}
  \mathlibok
  \uses{def:mathlib-triangulated-spectral-object}
  For a triangulated spectral object $T$ indexed by $I$ and a functor
  $G:I'\to I$, Mathlib constructs the triangulated spectral object
  $T\circ G$ indexed by $I'$.  Its functorial distinguished triangles are
  obtained by applying $G$ to composable pairs in $I'$.
\end{definition}

% SOURCE: Mathlib/Algebra/Homology/SpectralObject/Basic.lean, declaration CategoryTheory.Abelian.SpectralObject.
\begin{definition}[Mathlib abelian spectral object]
  \label{def:mathlib-spectral-object}
  \lean{CategoryTheory.Abelian.SpectralObject}
  \mathlibok
  A Mathlib spectral object in an abelian category is an indexed family of
  homology functors equipped with functorial connecting maps and the three
  exactness conditions forming its long exact sequences.
\end{definition}

% SOURCE: Mathlib/Algebra/Homology/SpectralObject/HasSpectralSequence.lean:67, declaration CategoryTheory.Abelian.SpectralObject.SpectralSequenceDataCore.
\begin{definition}[Mathlib spectral-sequence recipe]
  \label{def:mathlib-spectral-sequence-data-core}
  \lean{CategoryTheory.Abelian.SpectralObject.SpectralSequenceDataCore}
  \mathlibok
  \uses{def:mathlib-spectral-object,def:mathlib-complex-shape-up}
  A spectral-sequence data core records the grading degree and four monotone
  spectral-object indices that produce each page object and its differential
  shape.
\end{definition}

% SOURCE: Mathlib/Algebra/Homology/SpectralObject/HasSpectralSequence.lean:303, declaration CategoryTheory.Abelian.SpectralObject.HasSpectralSequence.
\begin{definition}[Mathlib spectral-object admissibility]
  \label{def:mathlib-has-spectral-sequence}
  \lean{CategoryTheory.Abelian.SpectralObject.HasSpectralSequence}
  \mathlibok
  \uses{def:mathlib-spectral-object,def:mathlib-spectral-sequence-data-core}
  For a spectral object and a page recipe, Mathlib's
  \texttt{HasSpectralSequence} class contains the endpoint vanishing needed to
  identify page homology with the successor page.
\end{definition}

% SOURCE: Mathlib/Algebra/Homology/SpectralObject/SpectralSequence.lean:481, declaration CategoryTheory.Abelian.SpectralObject.spectralSequence.
\begin{definition}[Mathlib spectral sequence of a spectral object]
  \label{def:mathlib-spectral-object-spectral-sequence}
  \lean{CategoryTheory.Abelian.SpectralObject.spectralSequence}
  \mathlibok
  \uses{def:mathlib-spectral-object,def:mathlib-spectral-sequence-data-core,def:mathlib-has-spectral-sequence,def:mathlib-spectral-sequence}
  Given a spectral object, a data core, and the endpoint-vanishing instance,
  Mathlib constructs a \texttt{CategoryTheory.SpectralSequence} with the
  prescribed page shapes.
\end{definition}

% SOURCE: Mathlib/Algebra/Homology/SpectralObject/SpectralSequence.lean:650, declaration CategoryTheory.Abelian.SpectralObject.E₂SpectralSequence.
\begin{definition}[Mathlib $E_2$ spectral sequence]
  \label{def:mathlib-e2-spectral-sequence}
  \lean{CategoryTheory.Abelian.SpectralObject.E₂SpectralSequence}
  \mathlibok
  \uses{def:mathlib-spectral-object-spectral-sequence}
  A spectral object indexed by extended integers has Mathlib's canonical
  cohomological $E_2$ spectral sequence on $\mathbb Z\times\mathbb Z$.
\end{definition}

% SOURCE: Mathlib/Algebra/Homology/SpectralObject/FirstPage.lean:78--84, declaration CategoryTheory.Abelian.SpectralObject.spectralSequenceFirstPageXIso.
\begin{definition}[Mathlib first-page computation]
  \label{def:mathlib-spectral-sequence-first-page}
  \lean{CategoryTheory.Abelian.SpectralObject.spectralSequenceFirstPageXIso}
  \mathlibok
  \uses{def:mathlib-spectral-object-spectral-sequence,def:mathlib-spectral-sequence-data-core,def:mathlib-has-spectral-sequence}
  If a spectral object has the required spectral-sequence instance and its
  data core has the first-page equalities required for Mathlib's computation,
  then the object at a bidegree of its initial page is canonically isomorphic
  to the corresponding value of the underlying spectral object's homology
  functor on the selected interval.  This computes the initial page object; it
  does not identify it with an associated-graded object of a filtered complex.
\end{definition}

% SOURCE: Mathlib/Algebra/Homology/SpectralObject/Cycles.lean:58, declaration CategoryTheory.Abelian.SpectralObject.cycles.
\begin{definition}[Mathlib spectral-object cycles]
  \label{def:mathlib-spectral-object-cycles}
  \lean{CategoryTheory.Abelian.SpectralObject.cycles}
  \mathlibok
  \uses{def:mathlib-spectral-object}
  Mathlib's cycle object for a spectral-object interval is the appropriate
  kernel of its connecting morphism.
\end{definition}

% SOURCE: Mathlib/Algebra/Homology/SpectralObject/Cycles.lean:62, declaration CategoryTheory.Abelian.SpectralObject.opcycles.
\begin{definition}[Mathlib spectral-object opcycles]
  \label{def:mathlib-spectral-object-opcycles}
  \lean{CategoryTheory.Abelian.SpectralObject.opcycles}
  \mathlibok
  \uses{def:mathlib-spectral-object}
  Mathlib's opposite-cycle object is the corresponding cokernel.  It is
  categorical infrastructure, not yet the paper's chosen element subgroup
  $B_r$.
\end{definition}

% SOURCE: Mathlib/Algebra/Homology/HomotopyCategory/SpectralObject.lean:47, declaration HomotopyCategory.spectralObjectMappingCone.
\begin{definition}[Mapping-cone spectral object]
  \label{def:mathlib-mapping-cone-spectral-object}
  \lean{HomotopyCategory.spectralObjectMappingCone}
  \mathlibok
  \uses{def:mathlib-triangulated-spectral-object,def:mathlib-cochain-complex}
  For a preadditive category $\mathcal C$ with a zero object and binary
  biproducts, Mathlib's declaration has type
  \[
    \operatorname{Triangulated.SpectralObject}
      \bigl(\operatorname{HomotopyCategory}_{\operatorname{up}(\mathbb Z)}
      (\mathcal C),\operatorname{CochainComplex}(\mathcal C,\mathbb Z)\bigr).
  \]
  Thus it is a triangulated spectral object valued in the homotopy category of
  cochain complexes.  It neither consumes a chain complex nor produces an
  abelian spectral object; those two conversions are project obligations below.
\end{definition}

\section{Pages, convergence, and comparison}

% SOURCE: Mathlib/CategoryTheory/Triangulated/SpectralObject.lean:20-24 records this bridge as TODO;
% Mathlib/CategoryTheory/Triangulated/HomologicalFunctor.lean:217-240 supplies the three exactness lemmas.
\begin{theorem}[Exactness obligations for a homological image]
  \label{thm:homological-spectral-object-exactness}
  \uses{def:mathlib-triangulated-spectral-object,def:mathlib-homological-functor,def:mathlib-homological-shift-sequence,def:mathlib-homology-sequence-connecting,thm:mathlib-homology-sequence-exact-one,thm:mathlib-homology-sequence-exact-two,thm:mathlib-homology-sequence-exact-three}
  \lean{KIP126.Core.SpectralSequence.homologicalImage_exactness}
  \leanok
  Let $T$ be a triangulated spectral object in a pretriangulated category
  $\mathcal C$, indexed by $I$, and let $F:\mathcal C\to\mathcal A$ be a
  homological functor to an abelian category, equipped with a $\mathbb Z$-shift
  sequence $(F_n)$.  For every composable pair $i\to j\to k$ in $I$, applying
  $(F_n)$ to the distinguished triangle of $T$ gives three exact short
  complexes: exactness at the middle term, exactness at the third term followed
  by the connecting map to $F_{n+1}$, and exactness at the shifted first term.
\end{theorem}

\begin{proof}
  \leanok
  In the construction, use Mathlib's homology-sequence connecting morphism on
  the distinguished triangle attached to $i\to j\to k$, then apply the mapped
  middle-exactness, connecting-map exactness, and shifted exactness theorems;
  the latter two declarations implement the required rotation and inverse
  rotation.  The Lean theorem projects the three exactness fields together.
\end{proof}

\begin{theorem}[Naturality of the homological-image connecting maps]
  \label{thm:homological-image-delta-naturality}
  \uses{def:mathlib-triangulated-spectral-object,def:mathlib-homological-functor,def:mathlib-homological-shift-sequence,thm:mathlib-homology-sequence-connecting-naturality}
  \lean{KIP126.Core.SpectralSequence.homologicalImage_delta_naturality}
  \leanok
  For every shift step $n_0+1=n_1$, the connecting transformation in the
  homological image is natural in morphisms of composable pairs in the index
  category.
\end{theorem}

\begin{proof}
  \leanok
  This is the naturality field of the connecting natural transformation,
  obtained from Mathlib's connecting-map naturality theorem after applying the
  shift sequence to the functorial distinguished triangles.
\end{proof}

% SOURCE: project bridge required by Mathlib/CategoryTheory/Triangulated/SpectralObject.lean:20-24.
\begin{definition}[Homological image as an abelian spectral object]
  \label{def:homological-image-abelian-spectral-object}
  \uses{def:mathlib-triangulated-spectral-object,def:mathlib-spectral-object,def:mathlib-homological-functor,def:mathlib-homological-shift-sequence,def:mathlib-homology-sequence-connecting,thm:mathlib-homology-sequence-connecting-naturality,thm:homological-spectral-object-exactness,thm:homological-image-delta-naturality}
  \lean{KIP126.Core.SpectralSequence.homologicalImage}
  \lean{KIP126.Core.SpectralSequence.homologicalImage_H}
  \lean{KIP126.Core.SpectralSequence.homologicalImage_delta_app}
  \leanok
  For $T$, $F$, and $(F_n)$ as above, define an abelian spectral object
  $\operatorname{HImage}(T,F)$ by
  \[
    H^n_T(i\to j):=F_n\bigl(T(i\to j)\bigr).
  \]
  Its connecting morphisms are the homology-sequence connecting maps of the
  functorial distinguished triangles of $T$, and its three exactness fields are
  the three obligations in
  Theorem~\ref{thm:homological-spectral-object-exactness}.  The connecting-map
  naturality and all three exactness fields are now implemented by the
  project-owned construction \texttt{homologicalImage}.
\end{definition}

% SOURCE: KIP126/Core/SpectralSequence/HomologicalImage.lean:158-194, homological-image morphism and functor.
\begin{definition}[Functoriality of the homological image]
  \label{def:homological-image-functoriality}
  \uses{def:homological-image-abelian-spectral-object,def:mathlib-triangulated-spectral-object,def:mathlib-homological-functor,def:mathlib-homological-shift-sequence}
  \lean{KIP126.Core.SpectralSequence.homologicalImageHom}
  \lean{KIP126.Core.SpectralSequence.homologicalImageFunctor}
  \leanok
  A morphism of triangulated spectral objects is sent, after applying a fixed
  homological functor and its shift sequence, to a morphism of abelian spectral
  objects.  Naturality of the connecting maps proves the commutative square;
  identities and composition are preserved, so the bridge is functorial.
\end{definition}

\begin{definition}[Mapping-cone homological image]
  \label{def:mapping-cone-homological-image}
  \uses{def:mathlib-mapping-cone-spectral-object,def:homological-image-abelian-spectral-object}
  \lean{KIP126.Core.SpectralSequence.mappingConeHomologicalImage}
  \lean{KIP126.Core.SpectralSequence.mappingConeHomologicalImage_H}
  \lean{KIP126.Core.SpectralSequence.mappingConeHomologicalImage_exactness}
  \leanok
  For a preadditive category with a zero object and binary biproducts, and a
  homological functor with its integer shift sequence, specialize the preceding
  bridge to Mathlib's mapping-cone spectral object.  The result is an abelian
  spectral object indexed by cochain complexes, with exactness inherited from
  the distinguished mapping-cone triangles.
\end{definition}

% SOURCE: KIP126/Core/Algebra/Completion.lean, quotient-tower and completion-witness interface.
\begin{definition}[Degreewise completion witness]
  \label{def:core-filtration-completion-witness}
  \uses{def:core-filtration,def:core-eventually-zero-filtration,def:core-bounded-above-filtration,def:mathlib-abelian}
  \lean{KIP126.Core.Algebra.Filtration.quotientAt}
  \lean{KIP126.Core.Algebra.Filtration.quotientProjection}
  \lean{KIP126.Core.Algebra.Filtration.quotientTransition}
  \lean{KIP126.Core.Algebra.Filtration.quotientProjection_transition}
  \lean{KIP126.Core.Algebra.Filtration.quotientTransition_comp}
  \lean{KIP126.Core.Algebra.Filtration.quotientTransition_id}
  \lean{KIP126.Core.Algebra.Filtration.quotientTower}
  \lean{KIP126.Core.Algebra.Filtration.quotientTowerCone}
  \lean{KIP126.Core.Algebra.Filtration.CompletionWitness}
  \lean{KIP126.Core.Algebra.Filtration.CompletionWitness.hasLimit}
  \lean{KIP126.Core.Algebra.Filtration.CompletionWitness.completionIso}
  \lean{KIP126.Core.Algebra.Filtration.CompletionWitness.completionIso_hom_comp_limit_π}
  \lean{KIP126.Core.Algebra.Filtration.eq_bot_of_le_of_eq_bot}
  \lean{KIP126.Core.Algebra.Filtration.quotientProjection_isIso_of_eq_bot}
  \lean{KIP126.Core.Algebra.Filtration.quotientTransition_isIso_of_eq_bot}
  \lean{KIP126.Core.Algebra.Filtration.quotientTower_isEventuallyConstantTo_of_eq_bot}
  \lean{KIP126.Core.Algebra.Filtration.quotientTowerCone_isLimit_of_eq_bot}
  \lean{KIP126.Core.Algebra.Filtration.CompletionWitness.of_isEventuallyZero}
  \lean{KIP126.Core.Algebra.Filtration.CompletionWitness.of_isBoundedAbove}
  \leanok
  For a filtered graded object $A$ and a grading degree $i$, the quotient
  tower has value $A_i/F^sA_i$ at $s$ and transition
  $A_i/F^sA_i\to A_i/F^tA_i$ when $t\le s$.  Its canonical cone has point
  $A_i$ and legs the quotient projections.  A degreewise completion witness
  is precisely a proof that this canonical cone is limiting, and it supplies
  existence of the corresponding limit.  With that limit chosen, it gives the
  canonical isomorphism
  \[
    A_i\ \cong\ \varprojlim_s A_i/F^sA_i,
  \]
  characterized by its composites with every quotient projection.  The
  witness is explicit data: this definition does not prove that arbitrary
  filtrations are complete, separated, exhaustive, or convergent.  It does
  prove that a degreewise eventually-zero filtration, and in particular a
  degreewise bounded-above filtration, supplies this completion witness.
\end{definition}

% SOURCE: aimpaper/main.tex:253-271 and 827-862; strong convergence is assumed throughout.
\begin{definition}[Complete and separated decreasing filtration]
  \label{def:complete-separated-filtration}
  \uses{def:core-filtration,def:core-filtration-completion-witness}
  \notready
  For a decreasing filtration $F$ on a graded abelian group $A$, ordinary
  exhaustiveness means $\bigcup_{s\in\mathbb Z}F^sA=A$, separatedness means
  $\bigcap_{s\in\mathbb Z}F^sA=0$, and complete means that the canonical map
  \[
    A\longrightarrow\varprojlim_s A/F^sA
  \]
  is an isomorphism, degree by degree.  The quotient-tower part of this
  condition is represented by the explicit witness in
  Definition~\ref{def:core-filtration-completion-witness}; the combined
  complete, separated, and convergence interface remains pending.  These
  ordinary union/intersection conditions must not be confused with the
  stronger degreewise eventual-top predicate
  \texttt{Filtration.IsExhaustive} or eventual-bottom predicate
  \texttt{Filtration.IsEventuallyZero} from the Core.  A convergence witness
  must also record the regularity or boundedness condition
  that makes only finitely many filtration pieces relevant in each total
  degree.
\end{definition}

% SOURCE: aimpaper/main.tex:134-150, 264-295, 799-823, and 1277-1317.
\begin{definition}[Spectral-sequence page-level convention]
  \label{def:ss-page-level}
  \notready
  A page-level convention records the first page, successor operation, page
  index used in the chosen cycle/boundary presentation, and the translation
  from any auxiliary quotient level to that page.  In the AIM Adams
  convention the displayed page has the predecessor offset
  $E_r\cong Z_{r-1}/B_{r-1}$ for $r\ge2$, while the corresponding synthetic
  quotient is $S/\lambda^{r-1}$.  These three roles of $r$, $r-1$, and the
  quotient exponent are recorded explicitly and are never identified
  definitionally.
\end{definition}

% SOURCE: aimpaper/main.tex:134-150 and 799-807, Adams differential degree.
% VERIFIED-IN: KIP126/Classical/Adams/Basic.lean, declarations Bidegree,
% classicalAdamsShift, classicalAdamsTarget, classicalAdamsShape, and
% classicalAdamsShape_two_rel.
\begin{definition}[Compiled classical Adams $E_2\to E_3$ shape slice]
  \label{def:classical-adams-e2-e3-shape-slice}
  \uses{def:mathlib-spectral-sequence,def:mathlib-complex-shape-up,def:core-f2-module-category,thm:mathlib-module-category-abelian}
  \lean{KIP126.Classical.Adams.ClassicalAdamsSpectralSequence}
  \lean{KIP126.Classical.Adams.classicalAdamsShape_two_rel}
  \lean{KIP126.Classical.Adams.ClassicalAdamsSS.e₂ToE₃}
  \leanok
  The compiled tracer bullet is a Mathlib spectral sequence of
  $\mathbb F_2$-modules indexed by $\mathbb Z\times\mathbb Z$, displayed from
  $E_2$, whose page-$r$ differential shape is
  $\operatorname{up}'(r,r-1)$.  In particular, its $d_2$ component sends
  $(s,t)$ to $(s+2,t+1)$, and Mathlib's page-passage isomorphism identifies
  the homology of this $E_2$ page with $E_3$.
\end{definition}

% SOURCE: aimpaper/main.tex:134-150 and 799-807, Adams differential degree.
\begin{definition}[Classical and synthetic Adams complex shape]
  \label{def:adams-complex-shape}
  \uses{def:mathlib-complex-shape-up,def:ss-page-level}
  \notready
  On the bidegree index $\mathbb Z\times\mathbb Z$, the page-$r$ Adams
  differential has additive degree $(r,r-1)$, hence is represented by
  Mathlib's fixed-degree shape $\operatorname{up}'(r,r-1)$.  The synthetic
  shape uses the same first two coordinates and fixes the weight coordinate.
\end{definition}

% SOURCE: aimpaper/main.tex:264-295 and 999-1007, extension differential degree.
\begin{definition}[Extension spectral-sequence complex shape]
  \label{def:ess-complex-shape}
  \uses{def:mathlib-complex-shape-up,def:ss-page-level}
  \notready
  On the paper's $(s,t)$ indexing, an extension differential $d_n^f$ has
  additive degree $(n,n)$.  The corresponding Mathlib complex shape is
  specified pagewise by $\operatorname{up}'(n,n)$, with any spectrum side or
  synthetic weight retained as an additional fixed index.
\end{definition}

% SOURCE: aimpaper/main.tex:134-150 and 253-295.
\begin{definition}[Adams bidegree convention]
  \label{def:adams-bidegree}
  \lean{KIP126.Classical.Adams.Bidegree}
  \lean{KIP126.Classical.Adams.classicalAdamsShift}
  \lean{KIP126.Classical.Adams.classicalAdamsTarget}
  \notready
  A classical Adams class has bidegree $(s,t)$, Adams filtration $s$, and stem
  $t-s$.  A classical differential has
  \[
    d_r:E_r^{s,t}\longrightarrow E_r^{s+r,t+r-1},
  \]
  hence lowers the stem by one.  In an extension spectral sequence the page
  differential is instead written with the paper's filtration shift
  $(n,n)$; every comparison must record which convention it uses.
\end{definition}

% SOURCE: aimpaper/main.tex:690-739, 799-823, and 1138-1150.
\begin{definition}[Synthetic tridegree and page-level conversion]
  \label{def:synthetic-tridegree}
  \uses{def:adams-bidegree}
  \lean{KIP126.Synthetic.SpectralSequence.Tridegree}
  \lean{KIP126.Synthetic.SpectralSequence.syntheticAdamsShift}
  \lean{KIP126.Synthetic.SpectralSequence.syntheticAdamsTarget}
  \lean{KIP126.Synthetic.SpectralSequence.syntheticAdamsShape}
  \lean{KIP126.Synthetic.SpectralSequence.nuDegree}
  \lean{KIP126.Synthetic.SpectralSequence.lambdaDegree}
  \lean{KIP126.Synthetic.SpectralSequence.lambdaTarget}
  \lean{KIP126.Synthetic.SpectralSequence.SyntheticAdamsSS.e₂ToE₃}
  \lean{KIP126.Synthetic.SpectralSequence.SyntheticAdamsSS.d₂}
  \notready
  A synthetic Adams class is indexed by $(s,t,w)$, where $s$ is Adams
  filtration, $t-s$ is the topological stem, and $w$ is synthetic weight.  A
  synthetic $d_r$ has degree $(r,r-1,0)$.  Multiplication by $\lambda$ has
  degree $(0,0,-1)$.  The classical page $E_r$ corresponds to the quotient
  level $S/\lambda^{r-1}$; this translation is data, not a definitional
  equality.
\end{definition}

% SOURCE: aimpaper/main.tex:827-862, definitions of Z_r and B_r.
\begin{definition}[Chosen cycles, boundaries, and pages]
  \label{def:chosen-page-presentation}
  \uses{def:mathlib-spectral-sequence,def:adams-bidegree,def:ss-page-level,def:adams-complex-shape}
  \notready
  An elementwise presentation includes its displayed-page reindexing.  In the
  AIM Adams convention it chooses nested subgroups of the initial page
  $E_2$, with $Z_1=E_2$, $B_1=0$ and
  $B_r^{s,t}\subseteq Z_r^{s,t}\subseteq E_2^{s,t}$, and, for every
  $r\ge2$, identifications
  \[
    E_r^{s,t}\cong Z_{r-1}^{s,t}/B_{r-1}^{s,t}.
  \]
  The displayed differential is represented before quotienting as
  \[
    d_r:Z_{r-1}^{s,t}\longrightarrow
    E_2^{s+r,t+r-1}/B_{r-1}^{s+r,t+r-1};
  \]
  its kernel inside $Z_{r-1}$ defines $Z_r$.  The subgroup $B_r$ is the
  inverse image in $E_2$ of the differential image in $E_2/B_{r-1}$, so it
  contains $B_{r-1}$.
  A filtered spectral sequence initialized at $E_0$ must first apply its
  stored page reindexing before using these Adams symbols.  The presentation
  also includes
  $Z_\infty=\bigcap_r Z_r$ and $B_\infty=\bigcup_r B_r$ when these operations
  exist in the chosen algebraic setting.
\end{definition}

% SOURCE: aimpaper/main.tex:134-151, 827-862, and throughout the appendix.
\begin{definition}[Cycle, boundary, and permanent cycle]
  \label{def:ss-cycle-boundary-permanent}
  \uses{def:chosen-page-presentation}
  \notready
  A representative in the initial $E_2$ page is an $r$-cycle when it lies in
  $Z_r$, an $r$-boundary when it lies in $B_r$, and a permanent cycle when it
  lies in $Z_\infty$.  Thus an element of the displayed $E_r$ page is
  represented by an $(r-1)$-cycle modulo $(r-1)$-boundaries, and surviving
  the displayed $d_r$ upgrades its representative from $Z_{r-1}$ to $Z_r$.
  A page class survives to $E_\infty$ when it has a permanent representative;
  it may still represent zero if that representative lies in $B_\infty$.
\end{definition}

% SOURCE: aimpaper/main.tex:253-263, standing convergence assumptions.
\begin{definition}[Strong convergence with detection]
  \label{def:strong-convergence-detection}
  \uses{def:core-filtration,def:complete-separated-filtration,def:chosen-page-presentation}
  \notready
  A strongly convergent elementwise spectral sequence for a graded abutment
  $A_*$ consists of a complete, exhaustive, degreewise separated filtration
  $F^sA_n$ and compatible isomorphisms
  $E_\infty^{s,s+n}\cong F^sA_n/F^{s+1}A_n$.  A class
  $x\in E_\infty^{s,s+n}$ detects $\alpha\in A_n$ when $\alpha\in F^sA_n$
  maps to $x$ in this quotient.
\end{definition}

\begin{lemma}[Projection equation for a restricted differential]
  \label{lem:filtered-associated-graded-differential-projection}
  \uses{def:core-associated-graded-differential,def:core-filtered-chain-complex,def:core-to-associated-graded}
  \lean{KIP126.Core.SpectralSequence.FilteredComplex.toAssociatedGraded_comp_associatedGradedDifferential}
  \leanok
  The quotient projection from $F^sC_k$ to
  $\operatorname{gr}^s_F(C_k)$ followed by the associated-graded differential
  equals the chosen restriction of $d_k$ followed by the quotient projection at
  degree $k-1$.
\end{lemma}

\begin{proof}
  \leanok
  Apply the cokernel universal property to the defining factorisation of the
  restricted differential.
\end{proof}

\begin{definition}[Associated-graded chain complex]
  \label{def:filtered-associated-graded-complex}
  \uses{def:core-associated-graded-object,def:core-associated-graded-differential,thm:core-associated-graded-differential-sq}
  \leanok
  \lean{KIP126.Core.SpectralSequence.FilteredComplex.associatedGradedComplex}
  For a filtered chain complex $(C,F)$ and fixed filtration degree $s$, the
  associated-graded chain complex has degree-$k$ object
  $\operatorname{gr}^s_F(C_k)$ and differential the compiled
  associated-graded differential.  In the displayed consecutive component the
  differential is followed by the canonical index transport induced by
  $(k+1)-1=k$.  The compiled square-zero theorem supplies the chain-complex
  relation; this node packages those existing components as a Mathlib chain
  complex rather than introducing a second differential.
\end{definition}

\begin{theorem}[Displayed differential of the associated-graded complex]
  \label{thm:filtered-associated-graded-complex-d}
  \uses{def:filtered-associated-graded-complex,def:core-associated-graded-differential}
  \lean{KIP126.Core.SpectralSequence.FilteredComplex.associatedGradedComplex_d}
  \leanok
  In consecutive chain degrees, the differential of the associated-graded
  complex is the induced differential followed by the canonical transport
  along $(k+1)-1=k$.
\end{theorem}

\begin{proof}
  \leanok
  Unfold the chain-complex constructor and simplify its allowed differential
  component.
\end{proof}

% SOURCE: KIP126/Core/SpectralSequence/FilteredComplex.lean:641-764, associated-graded maps and functoriality.
\begin{definition}[Associated-graded map of a filtered chain map]
  \label{def:filtered-associated-graded-map}
  \uses{def:filtered-chain-map,def:core-associated-graded-map}
  \lean{KIP126.Core.SpectralSequence.FilteredComplex.Morphism.associatedGradedMap}
  \leanok
  At each filtration and chain degree, a filtered chain map induces the
  associated-graded map obtained from its filtration-level factorisation.
\end{definition}

\begin{theorem}[Naturality of the associated-graded differential]
  \label{thm:filtered-associated-graded-map-comm}
  \uses{def:filtered-chain-map,def:filtered-associated-graded-map,def:core-associated-graded-map,def:core-associated-graded-differential,thm:filtered-chain-map-restricted-differential,lem:core-associated-graded-map-projection,lem:filtered-associated-graded-differential-projection}
  \lean{KIP126.Core.SpectralSequence.FilteredComplex.Morphism.associatedGradedMap_comm}
  \leanok
  For every filtration degree $s$ and chain degree $k$, the map induced by a
  filtered chain map $u$ satisfies
  \[
    \operatorname{gr}^s(u)_k\,;\,d^{\operatorname{gr},G}_{s,k}
      =d^{\operatorname{gr},F}_{s,k}\,;\,\operatorname{gr}^s(u)_{k-1}.
  \]
\end{theorem}

\begin{proof}
  \leanok
  Precompose with the quotient epimorphism from $F^sC_k$.  The universal
  property of the two cokernels replaces the two associated-graded maps by
  their filtration-level factorisations, and the preceding restriction
  theorem identifies the middle composites.  The corresponding projection
  equation for the source differential then finishes by cancellation.
\end{proof}

\begin{definition}[Associated-graded chain map]
  \label{def:filtered-associated-graded-chain-map}
  \uses{def:filtered-chain-map,def:filtered-associated-graded-complex,def:filtered-associated-graded-map,def:core-associated-graded-map}
  \leanok
  \lean{KIP126.Core.SpectralSequence.FilteredComplex.Morphism.associatedGradedChainMap}
  A filtered chain map $u:(C,F)\to(D,G)$ induces a chain map on every fixed
  filtration grade.  Its component in chain degree $k$ is the compiled map
  \[
    \operatorname{gr}^s_F(C_k)\longrightarrow\operatorname{gr}^s_G(D_k),
  \]
  and the chain-map equation is compatibility with the compiled
  associated-graded differentials.
\end{definition}

\begin{definition}[Functorial associated-graded chain maps]
  \label{def:filtered-associated-graded-functoriality}
  \uses{def:filtered-associated-graded-chain-map,def:filtered-associated-graded-map,def:core-filtered-morphism-functoriality}
  \lean{KIP126.Core.SpectralSequence.FilteredComplex.Morphism.associatedGradedMap_id}
  \lean{KIP126.Core.SpectralSequence.FilteredComplex.Morphism.associatedGradedMap_comp}
  \lean{KIP126.Core.SpectralSequence.FilteredComplex.Morphism.associatedGradedChainMap_id}
  \lean{KIP126.Core.SpectralSequence.FilteredComplex.Morphism.associatedGradedChainMap_comp}
  \lean{KIP126.Core.SpectralSequence.FilteredComplex.associatedGradedFunctor}
  \leanok
  The associated-graded construction is functorial in filtered chain maps:
  identities induce identity maps, composites induce composites, and the
  fixed-filtration chain maps assemble into
  \texttt{associatedGradedFunctor}.  These equations are part of the compiled
  interface and are used by the later spectral-object adapter.
\end{definition}

% SOURCE: KIP126/Core/SpectralSequence/Convergence.lean, explicit selected-page/abutment comparison witness.
\begin{definition}[Pointwise selected-page comparison with a bounded complete endpoint abutment]
  \label{def:page-abutment-comparison-witness}
  \uses{def:filtered-complex-endpoint-e2-spectral-sequence,def:core-filtration-completion-witness,def:core-bounded-filtration,def:core-associated-graded,lem:core-bounded-below-exhaustive,lem:core-bounded-above-eventually-zero}
  \lean{KIP126.Core.SpectralSequence.PageAbutmentComparisonWitness}
  \lean{KIP126.Core.SpectralSequence.PageAbutmentComparisonWitness.completion}
  \lean{KIP126.Core.SpectralSequence.PageAbutmentComparisonWitness.exhaustive}
  \lean{KIP126.Core.SpectralSequence.PageAbutmentComparisonWitness.eventuallyZero}
  \leanok
  For one endpoint extension and one chosen homological functor, this witness
  records its endpoint boundary data, a graded abutment explicitly isomorphic
  to the shifted homological image of the upper endpoint, a decreasing
  filtration on that abutment, and degreewise boundedness.  It derives a
  canonical completion witness in every grading degree from the bounded-above
  component.  It also chooses, for every bidegree $(p,q)$, a page at least $2$
  and an isomorphism from that page object to the indicated associated-graded
  piece of the endpoint abutment.  The boundedness field yields compiled proofs
  of the strong degreewise eventual-top predicate named
  \texttt{IsExhaustive} and of degreewise eventual-bottom.  These imply
  ordinary exhaustiveness and separatedness, but are not definitions of those
  weaker notions.

  This is intentionally a selected-page comparison, not a strong-convergence
  theorem: it supplies neither coherence of comparisons under page passage nor
  a canonical $E_\infty$ object, and it does not identify the abutment with
  filtered homology of the original filtered complex.
\end{definition}

% SOURCE: aimpaper/main.tex:461-631 and 1218-1235, maps with shifted degrees.
\begin{definition}[Reindexed spectral-sequence morphism]
  \label{def:reindexed-ss-morphism}
  \uses{def:mathlib-graded-comap,def:mathlib-complex-shape-up,def:mathlib-spectral-sequence-hom,def:adams-bidegree,def:synthetic-tridegree}
  \lean{KIP126.Comparison.ClassicalSynthetic.ReindexedSpectralSequenceMap}
  \lean{KIP126.Comparison.ClassicalSynthetic.forgetWeightComparison}
  \lean{KIP126.Comparison.ClassicalSynthetic.forgetWeight_differential_degree}
  \notready
  A reindexed morphism specifies a map of page index types, a map of grading
  indices, page maps, and proofs that the page maps commute with differentials
  after the stated degree translation and with page passage.  The interface
  must cover shifts $(s,t)\mapsto(s+k,t+k+d)$ and the forgetful map
  $(s,t,w)\mapsto(s,t)$ used in classical--synthetic comparison.
\end{definition}

% SOURCE: aimpaper/main.tex:138-150, 159-162, 1606-1753.
\begin{definition}[Multiplicative spectral sequence]
  \label{def:multiplicative-ss}
  \uses{def:chosen-page-presentation}
  \notready
  A multiplicative spectral sequence has unital associative products on its
  pages, compatible page maps, additive bidegrees, and a Leibniz rule for page
  differentials with the chosen sign convention.  Its convergence data also
  identifies the $E_\infty$ product with the associated-graded product on the
  abutment.  This interface is required before expressions such as $h_j^2$ or
  $h_0x$ are meaningful.
\end{definition}

% SOURCE: aimpaper/main.tex:253-263 and all later elementwise arguments.
\begin{proposition}[Categorical and elementwise presentations agree]
  \label{prop:categorical-elementwise-comparison}
  \uses{def:mathlib-spectral-sequence,def:mathlib-short-complex-exact,thm:mathlib-exact-iff-image-kernel,def:mathlib-spectral-object-cycles,def:mathlib-spectral-object-opcycles,def:chosen-page-presentation,thm:filtered-complex-to-spectral-sequence}
  \notready
  For a spectral sequence in graded $\mathbb F_2$-vector spaces constructed
  from a chosen exact couple or filtered complex, the Mathlib page objects,
  after the stored page reindexing, are naturally isomorphic in the AIM Adams
  convention to $Z_{r-1}/B_{r-1}$ on the displayed $E_r$ page.  The Mathlib
  page differential corresponds to the representative map
  $Z_{r-1}\to E_2/B_{r-1}$.
\end{proposition}

\begin{proof}
  Construct the isomorphism on $E_2=Z_1/B_1$, identify each next Mathlib page
  with homology, and induct on the displayed index $r$.  Passing from $E_r$ to
  $E_{r+1}$ takes the homology of
  $Z_{r-1}/B_{r-1}$ and gives precisely $Z_r/B_r$; naturality supplies
  compatibility with page passage.  The construction must be carried
  out for the concrete classical and synthetic presentations, not postulated
  for an arbitrary spectral sequence.
\end{proof}

\section{Filtered-complex realization}

% SOURCE: KIP126/Core/SpectralSequence/FilteredComplex.lean: filtered-complex morphism declarations; Mathlib/Algebra/Homology/HomotopyCategory/SpectralObject.lean.
\begin{definition}[Filtered chain map]
  \label{def:filtered-chain-map}
  \uses{def:mathlib-chain-complex,def:core-filtered-chain-complex,def:core-filtered-morphism}
  \leanok
  \lean{KIP126.Core.SpectralSequence.FilteredComplex.Morphism}
  \lean{KIP126.Core.SpectralSequence.FilteredComplex.Morphism.toFilteredMorphism}
  \lean{KIP126.Core.SpectralSequence.FilteredComplex.Morphism.id}
  \lean{KIP126.Core.SpectralSequence.FilteredComplex.Morphism.comp}
  For filtered chain complexes $(C,F)$ and $(D,G)$, a filtered chain map is a
  Mathlib chain-complex morphism $u:C\to D$ whose underlying graded morphism is
  filtration preserving.  Equivalently, for every $(s,k)$ it includes the
  restriction $F^sC_k\to G^sD_k$ commuting with the two subobject inclusions;
  the chain-map square makes these restrictions commute with the restricted
  differentials.  Identities and composites retain this data.
\end{definition}

% SOURCE: KIP126/Core/SpectralSequence/FilteredComplex.lean:193-344, fixed-level complexes and filtration inclusion chain maps.
\begin{definition}[Fixed-level filtered complexes and inclusions]
  \label{def:filtered-chain-level-complex-and-inclusions}
  \uses{def:core-filtered-chain-complex,def:core-filtration-inclusions}
  \lean{KIP126.Core.SpectralSequence.FilteredComplex.filteredChainComplex}
  \lean{KIP126.Core.SpectralSequence.FilteredComplex.filteredChainComplex_X}
  \lean{KIP126.Core.SpectralSequence.FilteredComplex.filteredChainComplex_d}
  \lean{KIP126.Core.SpectralSequence.FilteredComplex.filtration_inclusion_comp_differential}
  \lean{KIP126.Core.SpectralSequence.FilteredComplex.filtrationInclusionChainMap}
  \lean{KIP126.Core.SpectralSequence.FilteredComplex.filtrationInclusionChainMap_f}
  \lean{KIP126.Core.SpectralSequence.FilteredComplex.filtrationInclusionChainMap_id}
  \lean{KIP126.Core.SpectralSequence.FilteredComplex.filtrationInclusionChainMap_comp}
  \lean{KIP126.Core.SpectralSequence.FilteredComplex.instCategory}
  \leanok
  At a fixed filtration level, the restricted objects and differentials form a
  genuine chain complex.  Every order relation $t\le s$ gives a chain map
  $F^sC\to F^tC$; its component is the canonical filtration inclusion, and
  these maps satisfy identity and composition laws.  The explicit
  differential-compatibility lemma records why the inclusions are chain maps.
\end{definition}

\begin{definition}[Fixed-level filtered-complex functor]
  \label{def:filtered-chain-level-functor}
  \uses{def:filtered-chain-map,def:core-filtered-chain-complex,def:filtered-chain-level-complex-and-inclusions}
  \lean{KIP126.Core.SpectralSequence.FilteredComplex.Morphism.filteredChainMap}
  \lean{KIP126.Core.SpectralSequence.FilteredComplex.Morphism.filteredChainMap_id}
  \lean{KIP126.Core.SpectralSequence.FilteredComplex.Morphism.filteredChainMap_comp}
  \lean{KIP126.Core.SpectralSequence.FilteredComplex.filteredChainComplexFunctor}
  \lean{KIP126.Core.SpectralSequence.FilteredComplex.Morphism.filteredChainDiagramNatTrans}
  \leanok
  Restricting a filtered chain map to a fixed filtration degree gives a chain
  map between the two restricted complexes.  These restrictions preserve
  identities and composition, and assemble into a natural transformation of
  the decreasing filtration diagrams.
\end{definition}

\begin{theorem}[Compatibility of filtered restrictions]
  \label{thm:filtered-chain-map-restricted-differential}
  \uses{def:filtered-chain-map,def:core-filtered-chain-complex}
  \lean{KIP126.Core.SpectralSequence.FilteredComplex.preserves_differential}
  \leanok
  For a filtered chain map $u:(C,F)\to(D,G)$, write
  $\phi^u_{s,k}:F^sC_k\to G^sD_k$ for its chosen restriction and write
  $\delta^F_{s,k}$ and $\delta^G_{s,k}$ for the restricted differentials.
  Then the restriction squares commute:
  \[
    \phi^u_{s,k}\,;\,\delta^G_{s,k}
      =\delta^F_{s,k}\,;\,\phi^u_{s,k-1}.
  \]
\end{theorem}

\begin{proof}
  \leanok
  Postcompose both sides with the target filtration inclusion.  The two
  factorisation equations and the chain-map square identify the resulting
  composites; cancellation of the monomorphism given by that inclusion gives
  the displayed equality.
\end{proof}

\begin{definition}[Filtered complex as a cochain diagram]
  \label{def:filtered-complex-cochain-diagram}
  \uses{def:mathlib-chain-complex,def:mathlib-cochain-complex,def:core-filtered-chain-complex,def:core-filtration-inclusions,def:filtered-chain-level-complex-and-inclusions,def:filtered-chain-level-functor}
  \lean{KIP126.Core.SpectralSequence.FilteredComplex.filteredChainDiagram}
  \lean{KIP126.Core.SpectralSequence.FilteredComplex.filteredCochainComplex}
  \lean{KIP126.Core.SpectralSequence.FilteredComplex.filteredCochainDiagram}
  \lean{KIP126.Core.SpectralSequence.FilteredComplex.Morphism.filteredCochainDiagramNatTrans}
  \lean{KIP126.Core.SpectralSequence.FilteredComplex.filteredChainDiagramFunctor}
  \lean{KIP126.Core.SpectralSequence.FilteredComplex.filteredCochainDiagramFunctor}
  \leanok
  Let $\mathbb Z_{\mathrm{dec}}$ be the filtration category with a unique arrow
  $s\to t$ exactly when $t\le s$.  A filtered chain complex $(C,F)$ determines
  a functor
  \[
    D_F:\mathbb Z_{\mathrm{dec}}\longrightarrow
      \operatorname{CochainComplex}(\mathcal C,\mathbb Z),
    \qquad D_F(s)^n=F^sC_{-n},
  \]
  where $\mathcal C$ is the ambient abelian category.
  The differential is the restriction of $d_{-n}:C_{-n}\to C_{-n-1}$, and an
  arrow $s\to t$ is sent to the inclusion $F^sC\hookrightarrow F^tC$.  This
  records the reversal $k=-n$ and the decreasing-order convention explicitly;
  the implementation uses Mathlib's canonical chain/cochain equivalence rather
  than silently identifying the two abbreviations.  A filtered chain map
  induces a natural transformation of both the chain and cochain diagrams.
\end{definition}

\begin{definition}[Filtered-complex triangulated spectral-object adapter]
  \label{def:filtered-complex-triangulated-spectral-object}
  \uses{def:filtered-complex-cochain-diagram,def:mathlib-mapping-cone-spectral-object,def:mathlib-triangulated-spectral-object-precomp}
  \lean{KIP126.Core.SpectralSequence.cochainDiagramTriangulatedSpectralObject}
  \lean{KIP126.Core.SpectralSequence.filteredComplexTriangulatedSpectralObject}
  \lean{KIP126.Core.SpectralSequence.filteredComplexTriangulatedSpectralObject_ω₁}
  \leanok
  Precomposing Mathlib's mapping-cone triangulated spectral object with the
  cochain diagram of a filtered complex gives a triangulated spectral object
  indexed by the decreasing filtration category.  This adapter records the
  mapping-cone construction and the filtration-to-cochain index map without
  asserting any endpoint or convergence property.
\end{definition}

% SOURCE: KIP126/Core/SpectralSequence/SpectralObjectAdapter.lean:26-342, precomposition morphism and filtered triangulated adapter functor.
\begin{definition}[Morphisms and functor of the triangulated adapter]
  \label{def:filtered-complex-triangulated-spectral-object-morphism}
  \uses{def:filtered-complex-cochain-diagram,def:mathlib-mapping-cone-spectral-object,def:mathlib-triangulated-spectral-object-precomp}
  \lean{KIP126.Core.SpectralSequence.spectralObjectMapComposableArrows}
  \lean{KIP126.Core.SpectralSequence.spectralObjectArrowMap}
  \lean{KIP126.Core.SpectralSequence.spectralObjectMapComposableArrows_one_app}
  \lean{KIP126.Core.SpectralSequence.spectralObjectPrecompHom}
  \lean{KIP126.Core.SpectralSequence.cochainDiagramTriangulatedSpectralObjectHom}
  \lean{KIP126.Core.SpectralSequence.cochainDiagramTriangulatedSpectralObjectFunctor}
  \lean{KIP126.Core.SpectralSequence.filteredComplexTriangulatedSpectralObjectHom}
  \lean{KIP126.Core.SpectralSequence.filteredComplexTriangulatedSpectralObjectFunctor}
  \leanok
  A natural transformation between cochain diagrams induces a morphism of
  the precomposed triangulated spectral objects.  Naturality of the mapping-cone
  structure supplies the commutative square, and identity/composition laws
  make the construction a functor from cochain diagrams, hence from filtered
  complexes through their filtered cochain-diagram functor.
\end{definition}

\begin{definition}[Filtered-complex abelian spectral-object adapter]
  \label{def:filtered-complex-abelian-spectral-object}
  \uses{def:filtered-complex-triangulated-spectral-object,def:homological-image-abelian-spectral-object}
  \lean{KIP126.Core.SpectralSequence.cochainDiagramHomologicalImage}
  \lean{KIP126.Core.SpectralSequence.filteredComplexAbelianSpectralObject}
  \lean{KIP126.Core.SpectralSequence.cochainDiagramHomologicalImage_H}
  \lean{KIP126.Core.SpectralSequence.filteredComplexAbelianSpectralObject_H}
  \leanok
  Applying a homological functor with its shift sequence to the preceding
  triangulated adapter yields an abelian spectral object.  The functor and its
  shift sequence remain explicit parameters; the construction does not infer
  convergence or identify formal endpoint objects with integer filtration
  levels.
\end{definition}

% SOURCE: KIP126/Core/SpectralSequence/SpectralObjectAdapter.lean:359-394, homological-image adapter morphisms and functor.
\begin{definition}[Morphisms and functor of the abelian adapter]
  \label{def:filtered-complex-abelian-spectral-object-morphism}
  \uses{def:filtered-complex-triangulated-spectral-object-morphism,def:homological-image-abelian-spectral-object,def:homological-image-functoriality}
  \lean{KIP126.Core.SpectralSequence.cochainDiagramHomologicalImageHom}
  \lean{KIP126.Core.SpectralSequence.cochainDiagramHomologicalImageFunctor}
  \lean{KIP126.Core.SpectralSequence.filteredComplexAbelianSpectralObjectHom}
  \lean{KIP126.Core.SpectralSequence.filteredComplexAbelianSpectralObjectFunctor}
  \leanok
  Applying the homological-image bridge to the preceding morphisms gives
  morphisms of abelian spectral objects.  The resulting identity and
  composition equations are packaged as a functor, with the homological
  functor and its shift sequence remaining explicit parameters.
\end{definition}

% SOURCE: KIP126/Core/SpectralSequence/Convergence.lean, endpoint-extension and boundary-witness interface.
\begin{definition}[Endpoint extension of a filtered-complex diagram]
  \label{def:filtered-complex-endpoint-extension}
  \uses{def:filtered-complex-cochain-diagram,def:mathlib-cochain-complex}
  \lean{KIP126.Core.SpectralSequence.decreasingFiltrationEmbedding}
  \lean{KIP126.Core.SpectralSequence.EndpointExtension}
  \lean{KIP126.Core.SpectralSequence.EndpointExtension.finiteDiagram}
  \lean{KIP126.Core.SpectralSequence.EndpointExtension.bottomCone}
  \lean{KIP126.Core.SpectralSequence.EndpointExtension.topCocone}
  \lean{KIP126.Core.SpectralSequence.EndpointExtension.BoundaryWitness}
  \lean{KIP126.Core.SpectralSequence.EndpointExtension.finiteBottomCone}
  \lean{KIP126.Core.SpectralSequence.EndpointExtension.finiteTopCocone}
  \lean{KIP126.Core.SpectralSequence.EndpointExtension.finiteBottomIsLimit}
  \lean{KIP126.Core.SpectralSequence.EndpointExtension.finiteTopIsColimit}
  \leanok
  For the decreasing filtration diagram $D_F$, write
  $\iota:\mathbb Z_{\mathrm{dec}}\to\operatorname{EInt}$ for
  $s\mapsto-s$.  An endpoint extension is an
  $\operatorname{EInt}$-indexed cochain-complex diagram $D$ together with a
  natural isomorphism $D_F\cong D\circ\iota$.  Thus the two formal endpoint
  objects are kept distinct from every finite filtration level.

  The construction supplies the canonical cone from $D(\bot)$ and cocone to
  $D(\top)$ over the finite restriction.  A boundary witness is additional
  data saying that this cone is limiting and zero, that this cocone is
  colimiting, and that $D(\top)$ is isomorphic to a designated cochain
  complex.  The two universal properties are also transported across the
  finite comparison isomorphism to $D_F$.  This interface neither constructs
  an endpoint extension for every filtered complex nor asserts completeness,
  separatedness, or convergence of filtered homology.
\end{definition}

% SOURCE: KIP126/Core/SpectralSequence/Convergence.lean, endpoint-extended spectral-object and first-page construction.
\begin{definition}[Canonical $E_2$ sequence of an endpoint extension]
  \label{def:filtered-complex-endpoint-e2-spectral-sequence}
  \uses{def:filtered-complex-endpoint-extension,def:filtered-complex-triangulated-spectral-object,def:homological-image-abelian-spectral-object,def:mathlib-e2-spectral-sequence,def:mathlib-spectral-sequence-first-page}
  \lean{KIP126.Core.SpectralSequence.EndpointExtension.triangulatedSpectralObject}
  \lean{KIP126.Core.SpectralSequence.EndpointExtension.abelianSpectralObject}
  \lean{KIP126.Core.SpectralSequence.EndpointExtension.endpointAbutment}
  \lean{KIP126.Core.SpectralSequence.EndpointExtension.spectralSequence}
  \lean{KIP126.Core.SpectralSequence.EndpointExtension.firstPageXIso}
  \leanok
  Given an endpoint extension $D$ and a homological functor from the homotopy
  category of cochain complexes, equipped with its integer shift sequence,
  the mapping-cone and homological-image constructions yield an abelian
  spectral object indexed by $\operatorname{EInt}$.  Mathlib then supplies
  its canonical cohomological $E_2$ spectral sequence on
  $\mathbb Z\times\mathbb Z$.  At $(p,q)$, the compiled first-page
  isomorphism identifies the $E_2$ object with the degree-$(p+q)$ homology
  functor evaluated on the interval from $q$ to $q+1$ in that spectral object.
  It does not identify this page with the associated graded of filtered
  homology or establish stabilization.
\end{definition}

% SOURCE: KIP126/Core/SpectralSequence/Convergence.lean, endpoint-extension morphisms and induced spectral-object morphisms.
\begin{definition}[Functoriality of endpoint extensions]
  \label{def:filtered-complex-endpoint-extension-morphism}
  \uses{def:filtered-complex-endpoint-extension,def:filtered-complex-endpoint-e2-spectral-sequence,def:filtered-complex-triangulated-spectral-object-morphism,def:homological-image-functoriality}
  \lean{KIP126.Core.SpectralSequence.EndpointExtension.Hom}
  \lean{KIP126.Core.SpectralSequence.EndpointExtension.Hom.id}
  \lean{KIP126.Core.SpectralSequence.EndpointExtension.Hom.comp}
  \lean{KIP126.Core.SpectralSequence.EndpointExtension.Hom.triangulatedSpectralObjectHom}
  \lean{KIP126.Core.SpectralSequence.EndpointExtension.Hom.abelianSpectralObjectHom}
  \leanok
  Over a filtered chain map, a morphism of endpoint extensions is a natural
  transformation of the extended diagrams whose finite restriction agrees
  with the filtered cochain-diagram map through the two comparison
  isomorphisms.  Such morphisms have specified identities and composites.
  They induce morphisms of the associated triangulated spectral objects and,
  after a fixed homological functor is applied, of the resulting abelian
  spectral objects.  The present interface deliberately stops there: it does
  not require preservation of separately chosen boundary or page/abutment
  witnesses, and it does not claim a packaged Mathlib spectral-sequence
  morphism or pagewise naturality for the canonical $E_2$ sequences.
\end{definition}

% SOURCE: KIP126/Core/SpectralSequence/SpectralObjectAdapter.lean: project adapter declarations;
% Mathlib/Algebra/Homology/HomotopyCategory/SpectralObject.lean;
% Mathlib/CategoryTheory/Triangulated/SpectralObject.lean;
% Mathlib/Algebra/Homology/HomotopyCategory/HomologicalFunctor.lean;
% Mathlib/Algebra/Homology/HomotopyCategory/ShiftSequence.lean.
\begin{definition}[Filtered-complex to spectral-object adapter]
  \label{def:filtered-complex-spectral-object-adapter}
  \uses{def:filtered-complex-cochain-diagram,def:filtered-complex-triangulated-spectral-object,def:filtered-complex-triangulated-spectral-object-morphism,def:filtered-complex-abelian-spectral-object,def:filtered-complex-abelian-spectral-object-morphism,def:filtered-complex-endpoint-extension,def:filtered-complex-endpoint-e2-spectral-sequence,def:page-abutment-comparison-witness,def:filtered-complex-endpoint-extension-morphism,def:core-filtration-completion-witness,def:complete-separated-filtration,def:mathlib-mapping-cone-spectral-object,def:mathlib-triangulated-spectral-object-precomp,def:homological-image-abelian-spectral-object}
  \notready
  The core precomposition and homological-image steps are implemented by
  Definitions~\ref{def:filtered-complex-triangulated-spectral-object} and
  \ref{def:filtered-complex-abelian-spectral-object}; the conditional
  endpoint-extension and $E_2$ constructions are implemented by
  Definitions~\ref{def:filtered-complex-endpoint-extension} and
  \ref{def:filtered-complex-endpoint-e2-spectral-sequence}.  The full adapter
  from an arbitrary filtered complex remains pending.  A concrete application
  must supply an $\operatorname{EInt}$ extension of the cochain diagram and
  any required endpoint comparison witnesses; it must not identify either
  endpoint with an integer filtration level.  Apply Mathlib's
  triangulated-spectral-object precomposition construction to the mapping-cone
  spectral object and the extended diagram, producing a \emph{triangulated}
  spectral object.  Apply
  Definition~\ref{def:homological-image-abelian-spectral-object} to a chosen
  homological functor on the homotopy category of cochain complexes together
  with its shift sequence, to obtain the required \emph{abelian} spectral
  object.  A degree-zero homology specialization is a later explicit
  instance, not part of this generic adapter.  A
  degreewise \texttt{CompletionWitness} now expresses the canonical
  quotient-tower completion condition, and a selected-page comparison can be
  kept as explicit data by
  Definition~\ref{def:page-abutment-comparison-witness}.  Their existence,
  page-passage coherence, and comparison with filtered homology are still
  separate obligations.  Separatedness, the elementwise $Z/B$ presentation,
  and a functorial spectral-sequence morphism also remain unformalized here.
\end{definition}

% SOURCE: aimpaper/main.tex:253-271, Definition def:ess; standard filtered-complex construction used implicitly.
\begin{theorem}[Filtered complex produces a spectral sequence]
  \label{thm:filtered-complex-to-spectral-sequence}
  \uses{def:filtered-complex-spectral-object-adapter,def:filtered-complex-triangulated-spectral-object-morphism,def:filtered-complex-abelian-spectral-object-morphism,def:filtered-chain-map,def:filtered-associated-graded-chain-map,def:filtered-associated-graded-functoriality,def:core-associated-graded-object,def:core-associated-graded-differential,thm:core-associated-graded-differential-sq,def:mathlib-spectral-sequence-data-core,def:mathlib-has-spectral-sequence,def:mathlib-spectral-object-spectral-sequence,def:mathlib-spectral-object-cycles,def:mathlib-spectral-object-opcycles,def:chosen-page-presentation,def:complete-separated-filtration,def:ss-page-level}
  \notready
  A decreasingly filtered chain complex of $\mathbb F_2$-vector spaces whose
  differential preserves filtration has a functorial spectral sequence with
  $E_0$ the associated-graded complex.  Under exhaustive, complete,
  separated, and degreewise regular convergence data, its $E_\infty$ page is
  the associated graded of filtered homology.  After applying the stored
  displayed-page reindexing, the construction supplies the chosen
  $E_r\cong Z_{r-1}/B_{r-1}$ presentation.  Every filtered chain map induces a Mathlib
  spectral-sequence morphism; on $E_0$ it is the associated-graded chain map,
  and this assignment preserves identities and composition.
\end{theorem}

\begin{proof}
  Apply the adapter and use the compiled associated-graded differential and its
  square-zero theorem to identify the initial page as a complex.  Define the
  page recipe for the chosen $E_0$ convention,
  and discharge Mathlib's endpoint-vanishing
  \texttt{HasSpectralSequence} fields from regularity.  Invoke Mathlib's
  spectral-object construction rather than defining a second spectral
  sequence.  Identify its categorical cycles and opcycles with the filtration
  kernels and cokernels and hence, with the recorded predecessor offset, with
  $Z_{r-1}$ and $B_{r-1}$ on the displayed $E_r$ page.  Finally use
  completeness, separatedness, and degreewise regularity to eliminate the
  Boardman obstruction and identify the limiting filtration with filtered
  homology.  For a filtered chain map, apply functoriality of the cochain
  diagram, mapping-cone construction, and homological-image bridge; the initial
  component is the associated-graded chain map by construction.
\end{proof}
